\section{Executive Summary}

\subsection{What this paper establishes}

This is a bounded characterization of FL-BSA \ProductReleaseTag{}, not a generic product claim.
The paper, figures, and companion are tied to the annotated product tag object
\IdentityText{\ProductTagObjectRaw}, peeled product commit
\IdentityText{\ProductCommitRaw}, release-evidence workflow run \EvidenceRunRaw{} attempt
\EvidenceRunAttemptRaw{}, and whitepaper source commit \IdentityText{\WhitepaperCommitRaw}. The
primary intake ZIP has SHA-256 \IdentityText{\EvidenceBundleShaRaw}. Those identities prevent a
later product commit or a different evidence run from being silently described as v5.0.1.

The paper establishes four narrower propositions:

\begin{enumerate}
  \item the amplification branch preserves the generated fixture's gender selection-rate signal;
  \item the intrinsic branch's equal-rate post-label policy mechanically produces near-parity,
        showing that the two branch pathways remain separated;
  \item the exact release passed its configured multi-seed Gold scenario gate; and
  \item the release evidence is reproducible and internally linked at the byte, run, image, and
        certificate-hash levels described here.
\end{enumerate}

It does \emph{not} establish why a historical disparity arose, what would have happened to real
borrowers under another policy, or whether a lender complies with any law. Those questions require
deployment-specific data, legal analysis, model validation, and accountable human decisions.

\subsection{Findings at a glance}

\begin{KeyBox}
\begin{tabular}{@{}>{\bfseries}p{0.25\linewidth}p{0.66\linewidth}@{}}
Fairness signal & Historical AIR \HistoricalAIR{} and amplification AIR \AmplificationAIR{} are
both below the internally selected 0.80 point screen. The amplification 95\% interval is
[\AmplificationAIRLower{}, \AmplificationAIRUpper{}]. \\
Branch control & Intrinsic AIR \IntrinsicAIR{} [\IntrinsicAIRLower{}, \IntrinsicAIRUpper{}]
follows an equal-rate post-label policy. This is an expected mechanical result, not a causal
``debiasing'' conclusion; the displayed range is for format symmetry, not inference. \\
Quality & Overall quality is \AmplificationQuality{} for amplification and \IntrinsicQuality{}
for intrinsic. The producer's certificate utility field is unevaluated for both branches. \\
Robustness & \RobustnessTotalPasses{}/\RobustnessTotalRuns{} configured Gold scenario runs passed
across ten seeds and four scenarios. This is gate stability, not an interval for a population
parameter. \\
Utility study & Ten equal-size train-on-synthetic runs yield mean ROC AUC
\UtilitySyntheticAUCMean{} (range \UtilitySyntheticAUCMin{}--\UtilitySyntheticAUCMax{}) on a
held-out generated fixture, versus \UtilityBaselineAUC{} for the real-train fixture baseline.
Skill-normalised AUC retention averages \UtilitySkillRetentionMeanPct{}\%;
\UtilityBelowChanceSeedCount{} of ten runs are below chance, and train-on-synthetic predictive
utility is not established. \\
\end{tabular}
\end{KeyBox}

The point screen and confidence-interval relation are reported separately. A row is never labelled
``PASS'' or ``FAIL'' merely because its point estimate lies on one side of 0.80. The historical
interval crosses 0.80; the amplification interval is entirely below it; and the intrinsic range is
not treated as inferential because its approvals are policy-determined. Historical and
amplification p-values and intervals remain conditional on one generated fixture and the configured
sample size.

\subsection{Race reference policy}

The configuration retains \RaceConfiguredReference{} as the provenance reference for race. The
active multi-group comparison policy instead uses the group with the highest observed selection
rate, which is \RaceEffectiveReference{} in this run, consistent with the comparison convention in
the employment-selection guidance that inspired the internal heuristic \citep{eeoc_uniform_guidelines}.
Race values are preserved in the companion but are not plotted as headline results because the
smallest group has \RaceMinimumGroupN{} rows (\RaceMinimumGroupPct{}\%). It
\RaceCountFloorRelation{} the configured \RaceDisplayMinGroupN{}-row count condition and
\RaceShareFloorRelation{} the configured \RaceDisplayMinGroupPct{}\% share condition. This display
decision does not change the machine-readable calculation.

\subsection{Publication boundary}

The release disposition records:
\begin{itemize}
  \item \path{customer_evidence_eligible=false};
  \item \path{customer_evidence_disposition=characterization_only}; and
  \item \path{publication_status=candidate_not_published}.
\end{itemize}
The authoritative locations are
\path{intake/archive/v5.0.1-release-30765888408.json} and the Gold release disposition in the
companion; those fields are not in \path{intake/pack_intent.json}. The candidate PDF and companion
must remain paired. No public companion URL exists at this review stage.

A security-supersession hold also applies. The exact v5.0.1 dependency manifest pins
\texttt{cryptography} 49.0.0; CVE-2026-69247 covers versions from 44.0.0 up to but excluding
50.0.0, with the fix beginning at 50.0.0 \citep{github_cve_2026_69247}. Consequently v5.0.1 is
security-superseded and must remain unpublished as a current-release whitepaper. Before any
publication, the owner must choose whether to retain this as an explicitly archival v5.0.1
characterization or retire it in favour of a separately rebuilt successor paper. The annotated
v5.0.2 tag now exists at tag object
\IdentityText{3d27f7d17c2c853753d40cb883858617ead21677}, peeled product commit
\IdentityText{b246e39a23be938397a6d28612c776bedd8b42e2}, and successful release-evidence run
31087235319; its lock pins \texttt{cryptography} 50.0.0. That release supersedes v5.0.1. This paper
does not characterize the successor evidence and must never be relabelled or presented as a
current-release v5.0.2 paper.

This paper does not claim or authorize vendor-authored evidence, broad commercial authorization,
public Marketplace go-live, production live decisioning, or regulator approval. Those explicit
non-claims are document-wide publication boundaries, not optional wording.
